\documentclass[11pt]{article}
\usepackage{amsmath,amsthm,amssymb,mathtools}
\usepackage{cleveref}
\usepackage{minted}
\usepackage{titlesec}
\usepackage{changepage}
\usepackage[letterpaper,bottom=1in,left=1in,right=1in]{geometry}
\usepackage[textsize=scriptsize]{todonotes}

% \titleformat*{\section}{\normalsize\bfseries}
% \titleformat*{\subsection}{\small\bfseries}

\newcommand{\bb}[1]{\mathbb{#1}}

\renewcommand{\theequation}{\theproblem.\arabic{equation}}

% ------------------------------- PROBLEM COMMANDS ------------------------------- %
\newcounter{problem}
\setcounter{problem}{0}

\newcounter{subproblem}[problem]
\setcounter{subproblem}{0}
\renewcommand{\thesubproblem}{\theproblem\alph{subproblem}}

\newcommand{\problem}[1]{
    \refstepcounter{problem}
    \noindent\textbf{Prob \theproblem}. #1
}

\newcommand{\subproblem}[1]{
    \refstepcounter{subproblem}
    \vspace{1em}
    \noindent\hspace{2em}\theproblem\alph{subproblem}. #1
}
% ------------------------------- PROBLEM COMMANDS ------------------------------- %

% ------------------------------- NEW THEOREMS ------------------------------------ %
\newtheorem{lemma}{Lemma}[problem]

\newtheoremstyle{claim}% name
  {\topsep}% space above
  {\topsep}% space below
  {}% body font
  {}% indent amount
  {\itshape}% theorem head font
  {}% punctuation after theorem head
  {.5em}% space after theorem head
  {\thmname{#1}\thmnumber{ #2}\thmnote{ (#3)}}% theorem head spec
\theoremstyle{claim}
\newtheorem*{claim}{Claim.}

% ------------------------------- NEW THEOREMS ----------------------------------- %


% ------------------------------- SOLUTION COMMAND ------------------------------- %
\newcommand{\answer}[1]{
    \noindent Ans. #1
}
% ------------------------------- SOLUTION COMMAND ------------------------------- %


% ------------------------------- FANCY TITLE PAGE ------------------------------- %
\makeatletter
\newcommand{\@email}{jdoe@uni.edu}
\newcommand{\email}[1]{
    \renewcommand{\@email}{#1}
}

\newcommand{\@class}{CSCE 52003}
\newcommand{\class}[1]{
    \renewcommand{\@class}{#1}
}

\newcommand{\@classsection}{001}
\newcommand{\classsection}[1]{
    \renewcommand{\@classsection}{#1}
}

\newcommand{\@professor}{John Doe}
\newcommand{\professor}[1]{
    \renewcommand{\@professor}{#1}
}

\makeatletter
\newcommand{\@duedate}{\today}
\newcommand{\duedate}[1]{
    \renewcommand{\@duedate}{#1}
}

\renewcommand{\maketitle}{
    \begin{center}
        \normalsize{\@author \hfill \@duedate} \\
        \normalsize{\texttt{\@email} \hfill \@class\,-\,\@classsection\ --\ \@professor} \\
        \vspace{-0.75em}
        \rule{\textwidth}{0.05em} \\
        \vspace{0.75em}
        \Large\textbf{\@title} \\
        \vspace{1.5em}
    \end{center}
}
\makeatother
% ------------------------------- FANCY TITLE PAGE ------------------------------- %

\author{Brent Orlina}
\email{borlina@uark.edu} 
\class{CSCE 41303}
\classsection{001}
\professor{Trent Rogers}
\duedate{September 1, 2026}
\title{Problem Set 1}

\begin{document}
\maketitle

% --------- Problem 1 --------- %
\problem{%
    Indicate for each pair of expressions \((A, B)\) in the table below where \(A\) is \(O\), \(o\), \(\Omega\), \(\omega\), or \(\Theta\) of \(B\) (e.g. is \(\lg^k(n) \in O(n^\epsilon)?\). Assume that \(k \geq 1, \epsilon > 0\), and \(c > 1\) are constants. Write you answer in the form of the table with ``yes'' or ``no'' written in each box.

    \begin{center}
    \begin{tabular}{ cc|c|c|c|c|c| }
        \(A\)         & \(B\)              & \(O\) & \(o\) & \(\Omega\) & \(\omega\) & \(\Theta\) \\
        \hline
        \(\lg^k n\)   & \(n^\epsilon\)     &       &       &            &            & \\
        \hline
        \(n^k\)       & \(c^n\)            &       &       &            &            & \\
        \hline
        \(\sqrt{n}\)  & \(n^{\sin n}\)     &       &       &            &            & \\
        \hline
        \(2^n\)       & \(2^\frac{n}{2}\)  &       &       &            &            & \\
        \hline
        \(n^{\lg c}\) & \(c^{\lg n}\)      &       &       &            &            & \\
        \hline
        \(\lg(n!)\)   & \(\lg(n^n)\)        &       &       &            &            & \\
        \hline
    \end{tabular}
    \end{center}
}
% --------- Problem 1 --------- %

\pagebreak

% --------- Problem 2 --------- %
\problem{%
    Let \(f(n)\) and \(g(n)\) be asymptotically nonnegative functions. Show that \(f(n) \in o(g(n))\) if and only if \(g(n) \in \omega(f(n))\).
}
% --------- Problem 2 --------- %

\pagebreak

% --------- Problem 3 --------- %
\problem{%
    Prove or disprove \(2^n \in o(n!)\)
}
% --------- Problem 3 --------- %

\pagebreak

% --------- Problem 4 --------- %
\problem{%
    Use mathematical induction to show that when \(n \geq 2\) is an exact power of 2, the solution of the recurrence
    \[
        T(n) = \begin{cases}
            2                   &\text{if } n = 2 \\
            2T(\frac{n}{2}) + n &\text{if } n > 2
        \end{cases}
    \]
is \(T(n) = n \lg n\)
}
% --------- Problem 4 --------- %

\pagebreak

% --------- Problem 5 --------- %
\problem{%
    Prove the following algorithm correctly calculates the sum of an array \texttt{arr}.
}
\begin{minted}[linenos]{python}
def sum_array(arr, n):
    if n == 1:
        return arr[0]
    mid = n // 2
    left = sum_array(arr[:mid], mid)
    right = sum_array(arr[mid:], n - mid)
    return left + right
\end{minted}
% --------- Problem 5 --------- %

\end{document}

